\documentclass[11pt,a4paper]{article}

\usepackage[margin=1in]{geometry}
\usepackage{lmodern}
\usepackage{amsmath,amssymb,amsthm}
\usepackage{bm}
\usepackage{microtype}
\usepackage{hyperref}
\usepackage{xurl}
\usepackage{xcolor}
\hypersetup{colorlinks=true,linkcolor=blue,citecolor=blue,urlcolor=blue}
\sloppy

\newcommand{\id}{\mathbf{1}}          % identity matrix
\newcommand{\dd}{\mathrm{d}}
\newcommand{\ii}{\mathrm{i}}
\newcommand{\ee}{\mathrm{e}}
\newcommand{\tr}{\operatorname{tr}}
\newcommand{\Real}{\operatorname{Re}}
\newcommand{\Imag}{\operatorname{Im}}
\newcommand{\half}{\tfrac{1}{2}}
\newcommand{\gr}{\check{g}}            % bare lead GF (Keldysh-Nambu)
\newcommand{\Gr}{\check{G}}            % full GF
\newcommand{\Sig}{\check{\Sigma}}      % self-energy
\newcommand{\tauz}{\tau_3}
\newcommand{\Om}{\Omega}
\newcommand{\Wm}{W}

\title{\bf Keldysh--Floquet Formulation of Voltage- and
Current-Biased Josephson Junctions}
\author{Technical notes for the \texttt{Keldyshsetup\_Floquetn} codebase}
\date{}

\begin{document}
\maketitle

\begin{abstract}
These notes give a self-contained derivation of the equations solved by the
Julia codebase that accompanies the paper \emph{``Origin of Subharmonic Gap
Structure of DC Current-Biased Josephson Junctions''}
(\href{https://arxiv.org/abs/2412.09862}{arXiv:2412.09862}).  We set up the
nonequilibrium (Keldysh) Green's-function description of a single-channel
Josephson junction of arbitrary transparency in Nambu space, introduce the
Floquet (harmonic) representation generated by the time-dependent
superconducting phase, and write the Dyson, Keldysh and current equations.
We then specialize to the two driving protocols: \emph{voltage bias}
(constant $V$, a single-shot evaluation) and \emph{current bias} (constant
time-averaged voltage $\langle V\rangle$, requiring a self-consistent solution
for the phase dynamics).  Throughout we use units $\hbar=e=1$ and indicate the
correspondence to the routines in \texttt{Keldyshsetup\_Floquetn.jl}.
\end{abstract}

\tableofcontents

\section{Physical setup}
\label{sec:setup}

We model a Josephson junction as two BCS superconducting leads (left $L$ and
right $R$) coupled through a single conduction channel with hopping
(tunnelling) amplitude $\mathcal{T}$.  Each lead is a semi-infinite system
characterized by:
\begin{itemize}
  \item the superconducting gap $\Delta$ (all energies are measured in units of
        $\Delta$; the code sets \texttt{delta}\,$=1$),
  \item an internal hopping / bandwidth scale $\zeta$ (\texttt{zeta}), which
        controls the normal-state lead coupling and hence the normalization of
        the surface Green's function,
  \item a phenomenological Dynes-type broadening $\Gamma$ (\texttt{Gamma}) that
        regularizes the BCS singularities and supplies the lesser (occupation)
        self-energy through a coupling to equilibrium reservoirs.
\end{itemize}
The single dimensionless parameter that controls the junction transparency is
the ratio $\mathcal{T}/\zeta$: $\mathcal{T}/\zeta\ll 1$ is the tunnel limit,
while $\mathcal{T}/\zeta\to 1$ approaches the high-transparency (ballistic)
regime.  The chemical potential is $\mu=0$ (particle--hole symmetric point).

\subsection{Hamiltonian and the superconducting phase}

In Nambu (particle--hole) space we use the spinor
$\Psi_\alpha=(c_{\alpha\uparrow},\,c^\dagger_{\alpha\downarrow})^{\mathsf T}$,
$\alpha\in\{L,R\}$, and the Pauli matrices $\tau_0,\tau_1,\tau_2,\tau_3$ acting
on this space, with $\tauz=\mathrm{diag}(1,-1)$.  Each lead is described by a
mean-field BCS Hamiltonian with order parameter $\Delta\,\ee^{\ii\varphi_\alpha}$.
The junction Hamiltonian is
\begin{equation}
  H(t) = H_L + H_R + H_T(t),
\end{equation}
where the coupling term, after a gauge transformation that moves the entire
phase difference $\phi(t)=\varphi_L-\varphi_R$ onto the hopping, reads
\begin{equation}
  H_T(t) = \mathcal{T}\,
  \Psi_L^\dagger
  \begin{pmatrix}
     \ee^{-\ii\phi(t)/2} & 0\\[2pt]
     0 & -\,\ee^{+\ii\phi(t)/2}
  \end{pmatrix}
  \Psi_R + \text{h.c.}
  \label{eq:HT}
\end{equation}
The two Nambu entries carry opposite phase factors $\ee^{\mp\ii\phi/2}$ because
the electron and hole components couple to the pairing field with opposite
charge.

\subsection{Josephson relation and the two bias protocols}
\label{sec:bias}

The phase difference is tied to the voltage drop across the junction by the
Josephson (gauge) relation
\begin{equation}
  \dot\phi(t) = 2\,V(t).
  \label{eq:josephson}
\end{equation}
Two distinct steady-state driving protocols follow:

\paragraph{Voltage bias.}
$V(t)=V_{\rm dc}$ is held constant.  Then $\phi(t)=2V_{\rm dc}\,t$ and
\begin{equation}
  \ee^{-\ii\phi(t)/2} = \ee^{-\ii V_{\rm dc} t}
  \quad\Longrightarrow\quad
  \mathcal{T}\,\ee^{-\ii\phi/2}\ \text{is a single harmonic at}\ \Om\equiv V_{\rm dc}.
  \label{eq:vbias_phase}
\end{equation}
No internal degrees of freedom remain; the current is obtained by a single
evaluation of the Green's functions.  In the code this corresponds to fixing
the harmonic vector to a single nonzero entry,
\verb|VipI[2*Nf] = 1| (see \texttt{Josephson\_Vbias\_Floquetn.jl}).

\paragraph{Current bias (the central problem).}
A fixed dc current is imposed.  Only the \emph{time-averaged} voltage
$\langle V\rangle$ is fixed; the instantaneous voltage $V(t)$ --- equivalently
the phase $\phi(t)$ --- becomes a dynamical quantity that develops harmonics and
must be determined self-consistently.  Writing
\begin{equation}
  \phi(t) = 2\langle V\rangle\,t + \delta\phi(t),
  \qquad \delta\phi(t+T_J)=\delta\phi(t),
  \qquad T_J = \frac{2\pi}{\omega_J},\quad \omega_J = 2\langle V\rangle,
\end{equation}
the object that enters the coupling \eqref{eq:HT} is periodic up to the carrier
$\ee^{-\ii\langle V\rangle t}$.  Defining the fundamental Floquet frequency
\begin{equation}
  \boxed{\ \Om \equiv \langle V\rangle\ }
\end{equation}
(this is the variable \texttt{Omega}\,$=$\,\texttt{ev} in the code), the gauge
factor admits the Fourier expansion
\begin{equation}
  \ee^{-\ii\phi(t)/2} \;=\; \sum_{m\ \mathrm{odd}} \Wm_m\,\ee^{-\ii m\Om t},
  \label{eq:Wexpansion}
\end{equation}
with \emph{only odd harmonics} $m$ present (the carrier $m=1$ plus the even
harmonics of $\delta\phi/2$).  The complex coefficients $\{\Wm_m\}$ are the
fundamental unknowns of the current-bias problem.  Physical observables (the
voltage $V(t)$ and the current $I(t)$) are periodic at the Josephson frequency
$\omega_J=2\Om$, i.e.\ they contain only \emph{even} harmonics of $\Om$.

\section{Nonequilibrium Green's functions in Nambu space}
\label{sec:gf}

We work with the contour-ordered Green's functions in the Keldysh
representation.  For each lead and for the full junction we use the
$2\times2$ Nambu retarded ($r$), advanced ($a$) and lesser ($<$) blocks; the
advanced block is obtained from the retarded one by Hermitian conjugation,
$g^a=(g^r)^\dagger$.

\subsection{Surface Green's function of a lead}
\label{sec:surface}

The local (surface) retarded Green's function of a semi-infinite BCS lead used
throughout the code is
\begin{equation}
  g^r(\omega) \;=\;
  \frac{1}{\zeta\,\sqrt{\Delta^2-(\omega+\ii\Gamma)^2}}
  \begin{pmatrix}
    -(\omega+\ii\Gamma) & \Delta\\[2pt]
    \Delta & -(\omega+\ii\Gamma)
  \end{pmatrix}
  \;=\;
  \frac{-(\omega+\ii\Gamma)\,\tau_0+\Delta\,\tau_1}
       {\zeta\,\sqrt{\Delta^2-(\omega+\ii\Gamma)^2}} ,
  \label{eq:surface_gf}
\end{equation}
where the branch of the square root is chosen so that $\Imag g^r\le 0$.  This is
implemented in \texttt{grwmnf} (retarded) and \texttt{gawmnf} (advanced); an
alternative transfer-matrix construction of the same surface Green's function
for a tight-binding chain (following Phys.\ Rev.\ B \textbf{83}, 085412 (2011))
is provided in \texttt{surfacegr}.

The bare lesser function follows from the fluctuation--dissipation relation at
zero temperature ($f(\omega)=\theta(-\omega)$):
\begin{equation}
  g^<(\omega) = -\,\bigl[g^r(\omega)-g^a(\omega)\bigr]\,\theta(-\omega)
              = 2\ii\,\theta(-\omega)\,\Imag\, g^r(\omega) ,
  \label{eq:bare_lesser}
\end{equation}
implemented in \texttt{glwmnf} (the step function $\theta(-\omega)$ is applied
mode by mode in the Floquet representation below).

\subsection{Floquet representation}
\label{sec:floquet}

Because the coupling \eqref{eq:HT} is periodic in time with fundamental
frequency $\Om$, any two-time function $A(t,t')$ that is periodic in the
center-of-mass time can be written in the Floquet (double-Fourier)
representation
\begin{equation}
  A(t,t') = \sum_{m,n}\int\frac{\dd\omega}{2\pi}\,
            \ee^{-\ii(\omega+m\Om)t}\,A_{mn}(\omega)\,\ee^{+\ii(\omega+n\Om)t'},
\end{equation}
so that products of time-local operators become matrix products over the
Floquet indices $m,n\in\{-N_f,\dots,N_f\}$ (the truncation \texttt{Nf}).  A
diagonal (time-local) bare propagator carries shifted energies,
\begin{equation}
  \bigl[g^r(\omega)\bigr]_{mn} = \delta_{mn}\,g^r(\omega+m\Om),
  \label{eq:floquet_bare}
\end{equation}
which is exactly the block-diagonal structure built by \texttt{grwmnf}.

The full working space is therefore
\begin{equation}
  \underbrace{2}_{L,R}\times
  \underbrace{2}_{\text{Nambu}}\times
  \underbrace{(2N_f+1)}_{\text{Floquet}}
  \;=\; 4(2N_f+1)
\end{equation}
dimensional.  Lead-diagonal quantities are block-diagonal in the $L/R$ index,
while the hopping self-energy is purely off-diagonal in $L/R$.

\subsection{Hopping self-energy in Floquet space}
\label{sec:hop}

Inserting the expansion \eqref{eq:Wexpansion} into the coupling \eqref{eq:HT},
the retarded hopping self-energy becomes a \emph{static} matrix in the Floquet
basis that couples the leads.  With the Floquet indices $m,n$ labelling rows and
columns (code block position $-m+N_f+1$), its lead-off-diagonal Nambu blocks are
\begin{equation}
  \bigl[\check V\bigr]^{LR}_{mn} =
  \mathcal{T}
  \begin{pmatrix}
    \Wm_{\,n-m} & 0\\[2pt]
    0 & -\,\Wm^{*}_{\,m-n}
  \end{pmatrix},
  \qquad
  \bigl[\check V\bigr]^{RL}_{mn} =
  \mathcal{T}
  \begin{pmatrix}
    \Wm^{*}_{\,m-n} & 0\\[2pt]
    0 & -\,\Wm_{\,n-m}
  \end{pmatrix},
  \qquad
  \check V^{LL}=\check V^{RR}=0 .
  \label{eq:hop_floquet}
\end{equation}
Here $\Wm_r$ are the Fourier coefficients of $\ee^{-\ii\phi/2}$ and
$\Wm^*_{r}$ those of $\ee^{+\ii\phi/2}$ (because $\phi$ is real), packed in the
code into the complex array \texttt{Vip} of length $4N_f+1$ whose entries are
nonzero only on odd harmonics.  This is assembled in \texttt{Vwmnf}.  The current
vertex carries the extra charge matrix $\tauz$ from the current operator; it is
identical to \eqref{eq:hop_floquet} but with the sign of the lower (hole) Nambu
entry flipped,
\begin{equation}
  \begin{aligned}
  \check V_i &\;=\; \big(\id_{LR}\otimes\id_{\rm Fl}\otimes\tauz\big)\,\check V ,\\[4pt]
  \bigl[\check V_i\bigr]^{LR}_{mn} &= \mathcal{T}\,\mathrm{diag}\!\big(\Wm_{\,n-m},\,+\Wm^{*}_{\,m-n}\big),
  \qquad
  \bigl[\check V_i\bigr]^{RL}_{mn} = \mathcal{T}\,\mathrm{diag}\!\big(\Wm^{*}_{\,m-n},\,+\Wm_{\,n-m}\big),
  \end{aligned}
  \label{eq:hop_vertex}
\end{equation}
built by \texttt{Viwmnf}.  Both $\check V$ and $\check V_i$ are \emph{linear} in
the unknowns $\{\Wm_r\}$ --- a fact that is central to the Jacobian construction
of \S\ref{sec:jacobian}.

\section{Dyson, Keldysh and current equations}
\label{sec:dyson}

\subsection{Dyson equation (retarded / advanced)}

The full retarded Green's function of the coupled junction follows from the
Dyson equation in the $4(2N_f+1)$-dimensional Floquet--Nambu--lead space,
\begin{equation}
  \Gr^r = \bigl(\,\id - \gr^r\,\Sig^r\,\bigr)^{-1}\gr^r,
  \qquad
  \Gr^a = (\Gr^r)^\dagger,
  \label{eq:dyson}
\end{equation}
where $\gr^r=\mathrm{diag}(g^r_L,g^r_R)$ is the block-diagonal bare propagator
\eqref{eq:floquet_bare} and $\Sig^r$ is the hopping self-energy
\eqref{eq:hop_floquet}.  This is \texttt{Grwmnf} (solved as a linear system,
not by explicit inversion).

\subsection{Lesser self-energy}
\label{sec:lesser_se}

The lesser self-energy collects the occupation information.  The code supports
two additive contributions (flags \texttt{Sigl\_b}, \texttt{Sigl\_s} in
\texttt{Siglf}):
\begin{align}
  \text{(bath)}\quad
  \Sigma^<_{\rm bath}(\omega+m\Om)
    &= 2\ii\,\Gamma\,\theta\!\left(-(\omega+m\Om)\right)\,\tau_0,
  \label{eq:sigl_bath}\\[4pt]
  \text{(surface)}\quad
  \Sigma^<_{\rm surf}(\omega+m\Om)
    &= \zeta^2\,\tauz\, g^<(\omega+m\Om)\,\tauz,
  \label{eq:sigl_surf}
\end{align}
applied mode by mode.  Equation \eqref{eq:sigl_bath} represents local
equilibrium reservoirs of strength $\Gamma$ attached to each site, while
\eqref{eq:sigl_surf} embeds the semi-infinite-lead occupation into the contact
site.

\subsection{Keldysh equation (lesser Green's function)}

The full lesser Green's function is obtained from the Keldysh equation
\begin{equation}
  \Gr^< = \Gr^r\,\Sig^<\,\Gr^a,
  \label{eq:keldysh}
\end{equation}
implemented in \texttt{Glesser\_Floquet\_Tfull}.

\subsection{Current formula}
\label{sec:current}

The charge current through the junction is the rate of change of charge on one
lead, which in terms of the contact self-energy and the lesser propagator reads
\begin{equation}
  I(t) = \Real\,\tr\!\Big[\tauz\big(\Sigma_{LR}\,G^<_{RL}-\Sigma_{RL}\,G^<_{LR}\big)\Big],
  \label{eq:current_time}
\end{equation}
where the trace is over Nambu space.  In the Floquet representation this becomes
a matrix in the Floquet indices, and the harmonic content of the current is
\begin{equation}
  I(t) = \sum_{h} I_{2h}\,\ee^{-\ii\,2h\Om t},
  \qquad
  I_{2h} = \frac{1}{2\pi}\int_{-\Om/2}^{+\Om/2}\!\dd\omega\;
           \big[\,\hat{I}(\omega)\,\big]^{\text{harmonic } 2h},
  \label{eq:current_harmonics}
\end{equation}
where only even harmonics $2h$ of $\Om$ (i.e.\ harmonics of the Josephson
frequency $\omega_J=2\Om$) appear.  The frequency integral runs over one
Floquet Brillouin zone $[-\Om/2,\Om/2)$ discretized with spacing
\texttt{dw0}; the energy grid is \texttt{war0}.  The dc current is the
$h=0$ component,
\begin{equation}
  I_{\rm dc} = I_0 = \frac{1}{2\pi}\int_{-\Om/2}^{+\Om/2}\!\dd\omega\;
               \tr_{\rm Nambu,Floquet}\!\big[\hat{I}(\omega)\big]_{\rm diagonal}.
\end{equation}
This is evaluated by \texttt{current\_Floquet\_Tfull} (the Floquet-mode-resolved matrix
\texttt{Iif}) and summed on the diagonal in \texttt{phisolve}.  The combination
$-\,\check V_i\,\Gr^<$ in \texttt{current\_Floquet\_Tfull} is exactly
\eqref{eq:current_time} written with the $\tauz$-dressed coupling
\eqref{eq:hop_floquet} from \texttt{Viwmnf}.

\section{Voltage-bias solution}
\label{sec:vbias}

For a constant voltage $V_{\rm dc}$, the phase gauge factor
\eqref{eq:vbias_phase} contains a single harmonic, so the coefficient vector
reduces to $\Wm_m=\delta_{m,1}$.  No self-consistency is required: one simply

\begin{enumerate}
  \item builds the bare propagators $g^{r,a,<}$ on the grid \texttt{war0};
  \item builds the (now trivial) hopping self-energy \eqref{eq:hop_floquet};
  \item solves the Dyson equation \eqref{eq:dyson} and Keldysh equation
        \eqref{eq:keldysh};
  \item evaluates the current \eqref{eq:current_time}.
\end{enumerate}
The dc component $I_{\rm dc}(V_{\rm dc})$ traced out as a function of bias gives
the voltage-biased $I$--$V$ curve.  Driver:
\texttt{Josephson\_Vbias\_Floquetn.jl}.  (A complementary direct
real-frequency / multiple-Andreev-reflection evaluation is provided by
\texttt{Josephson\_Vbias\_Floquetn\_direct.jl}.)

\section{Current-bias self-consistency}
\label{sec:ibias}

Under current bias the harmonic coefficients $\{\Wm_m\}$ are unknown and are
fixed by three sets of conditions.  Let the unknown real vector be the real and
imaginary parts of the $2N_f$ odd-harmonic coefficients (so $4N_f$ real
unknowns, packed in \texttt{Vipi}/\texttt{Vipsol}).

\paragraph{(i) Unitarity of the gauge factor.}
Since $\ee^{-\ii\phi/2}\,\ee^{+\ii\phi/2}=1$ identically, the convolution of the
coefficients of $\ee^{-\ii\phi/2}$ with those of $\ee^{+\ii\phi/2}$ must equal
the unit (zeroth) harmonic and vanish on all others:
\begin{equation}
  \sum_{k}\Wm_{k}\,\Wm^{*}_{k-2h} = \delta_{h,0},
  \qquad h=0,\pm1,\dots
  \label{eq:unitarity}
\end{equation}
This yields the real and imaginary constraint rows assembled at the top of
\texttt{IbiasResidual\_Tfull} (the products $\Real\Wm\,\Real\Wm+\Imag\Wm\,\Imag\Wm$ etc.),
including the normalization $\sum_k|\Wm_k|^2=1$.

\paragraph{(ii) Gauge / phase origin.}
The phase is pinned by $\phi(t{=}0)=0$, i.e.
\begin{equation}
  \sum_m \Imag\,\Wm_m = 0,
  \label{eq:gauge_fix}
\end{equation}
(the row \texttt{eqns[2*Nf] = sum(imag.(Vip))}).

\paragraph{(iii) Vanishing of all ac current harmonics (the bias condition).}
A pure dc current bias means \emph{no} alternating current flows; every nonzero
harmonic of the current \eqref{eq:current_harmonics} must vanish:
\begin{equation}
  \boxed{\;I_{2h} = 0 \quad\text{for all } h=1,\dots,N_f.\;}
  \label{eq:noac}
\end{equation}
Splitting into real and imaginary parts gives $2N_f$ equations, assembled in
\texttt{IbiasResidual\_Tfull} as
\begin{quote}
\ttfamily
eqns[2Nf+h] = Re(Ifa[-2h]), \quad eqns[2Nf+Nf+h] = Im(Ifa[-2h]).
\end{quote}

Conditions \eqref{eq:unitarity}--\eqref{eq:noac} together form a square
nonlinear system $\mathbf{F}(\mathbf{x})=\mathbf{0}$ of dimension $4N_f$, where
the current harmonics $I_{2h}$ are nonlinear functionals of $\{\Wm_m\}$ through
the Dyson/Keldysh equations \eqref{eq:dyson}--\eqref{eq:keldysh}.

\subsection{Numerical solution}

The system is solved with a Newton/trust-region method (\texttt{NLsolve.nlsolve}
with \texttt{method=:trust\_region}) using the analytically constructed
Jacobian (\texttt{IbiasJacobian\_Tfull}).  The bias is swept from high to low
$\langle V\rangle$, and the converged coefficients at one bias point seed the
solve at the next (numerical continuation).  Driver routine: \texttt{phisolve}
in \texttt{Keldyshsetup\_Floquetn.jl}; user driver:
\texttt{Josephson\_Ibias\_Floquetn.jl}.

Once $\{\Wm_m\}$ are known, the physical voltage waveform is reconstructed from
$\phi(t)=2\ii\log\!\big(\ee^{-\ii\phi(t)/2}\big)$ and $V(t)=\half\dot\phi$
(routine \texttt{Vt}), and the dc current $I_{\rm dc}(\langle V\rangle)$ is the
current-biased $I$--$V$ characteristic.

\section{Explicit Floquet-component form of the equations}
\label{sec:explicit}

We now write out, component by component, the nonlinear system
$\mathbf F(\mathbf x)=\mathbf 0$ that \texttt{IbiasResidual\_Tfull} assembles and
\texttt{phisolve} solves.  This is the precise content of \S\ref{sec:ibias}.

\subsection{The unknown vector}

Only odd harmonics of $\Om$ appear in $\ee^{-\ii\phi/2}$
\eqref{eq:Wexpansion}.  Index them by $q=1,\dots,2N_f$, corresponding to the odd
harmonics
\begin{equation}
  r_q = 2q-(2N_f+1)\in\{-(2N_f-1),\dots,-1,+1,\dots,+(2N_f-1)\}.
\end{equation}
The real unknown vector (the array \texttt{Vipi}, length $4N_f$) is
\begin{equation}
  \mathbf x = \big(\,\underbrace{\Real \Wm_{r_1},\dots,\Real \Wm_{r_{2N_f}}}_{2N_f},\;
                     \underbrace{\Imag \Wm_{r_1},\dots,\Imag \Wm_{r_{2N_f}}}_{2N_f}\,\big),
  \qquad
  \Wm_{r_q}=x_q+\ii\,x_{2N_f+q}.
  \label{eq:unknowns}
\end{equation}

\subsection{Current in Floquet components}
\label{sec:current_floquet}

Define the current matrix in the full lead$\otimes$Floquet$\otimes$Nambu space at
energy $\omega$,
\begin{equation}
  \hat I(\omega) \;\equiv\; -\,\check V_i(\Wm)\,\check G^<(\omega;\Wm),
  \label{eq:Ihat}
\end{equation}
which is the function \texttt{Itemp0}$=-$\texttt{Vvi}$*$\texttt{Glwmn}.  Its
Floquet-block- and Nambu-resolved trace, taking the $LL$ block minus the $RR$
block, gives the partial current carried between Floquet modes $m$ and $n$:
\begin{equation}
  I_{mn}(\omega) =
  \tr_{\rm Nambu}\!\Big[\hat I(\omega)^{LL}_{mn}\Big]
  -\tr_{\rm Nambu}\!\Big[\hat I(\omega)^{RR}_{mn}\Big].
  \label{eq:Imn_omega}
\end{equation}
Integrating over one Floquet zone $\omega\in[-\Om/2,\Om/2)$,
\begin{equation}
  I_{mn} = \frac{1}{2\pi}\int_{-\Om/2}^{\Om/2}\!\dd\omega\;I_{mn}(\omega)
         \;\approx\; \frac{\Delta\omega}{2\pi}\sum_{\omega\in\,\texttt{war0}} I_{mn}(\omega),
  \label{eq:Imn}
\end{equation}
the array \texttt{Iif}.  The $s$-th harmonic of the current
($I(t)=\sum_s I_s\,\ee^{-\ii s\Om t}$) is the sum over Floquet pairs with fixed
index difference,
\begin{equation}
  \boxed{\;I_s \;=\; \sum_{\substack{m,n=-N_f\\ n-m=s}}^{N_f} I_{mn}\;}
  \qquad\Longleftrightarrow\qquad
  \texttt{Ifa}[-(m-n)+(2N_f{+}1)] \mathrel{+}= \texttt{Iif}[\dots],
  \label{eq:Is}
\end{equation}
the array \texttt{Ifa}.  Because the coupling has only odd harmonics, the current
carries only \emph{even} $s$.  The dc current is the diagonal sum
\begin{equation}
  I_{\rm dc} = I_0 = \sum_{m=-N_f}^{N_f} I_{mm} = \Real\,\tr\,\texttt{Iif}.
\end{equation}

\subsection{The residual equations $\mathbf F(\mathbf x)$}

The $4N_f$ scalar residuals returned by \texttt{IbiasResidual\_Tfull} are, in order:

\paragraph{Rows $1\ldots 2N_f-1$ --- unitarity/normalization of the gauge factor.}
Since $\ee^{-\ii\phi/2}\ee^{+\ii\phi/2}\equiv1$, the harmonic-$2h$ component of the
product must vanish for $h\neq0$ and equal unity for $h=0$.  Defining the
convolution
\begin{equation}
  C_{2h} \;\equiv\; \sum_{r} \Wm_{r+2h}\,\Wm^{*}_{r},
  \label{eq:Cdef}
\end{equation}
the constraints are
\begin{equation}
  \boxed{\;C_{2h}=\delta_{h,0}\;}\qquad h=0,1,\dots,N_f-1,
  \label{eq:unitarity_explicit}
\end{equation}
split into
\begin{align}
  F_{h}        &= \Real C_{2h} = \sum_r\big(\Real \Wm_{r+2h}\Real \Wm_r+\Imag \Wm_{r+2h}\Imag \Wm_r\big),
                 & h&=1,\dots,N_f-1, \label{eq:F_reC}\\
  F_{N_f-1+h}  &= \Imag C_{2h} = \sum_r\big(\Imag \Wm_{r+2h}\Real \Wm_r-\Real \Wm_{r+2h}\Imag \Wm_r\big),
                 & h&=1,\dots,N_f-1, \label{eq:F_imC}\\
  F_{2N_f-1}   &= C_0-1 = \sum_r|\Wm_r|^2-1. \label{eq:F_norm}
\end{align}

\paragraph{Row $2N_f$ --- phase origin.}
The gauge $\phi(t{=}0)=0$ makes $\ee^{-\ii\phi(0)/2}=\sum_r \Wm_r$ real, hence
\begin{equation}
  F_{2N_f} = \sum_r \Imag \Wm_r = 0.
  \label{eq:F_gauge}
\end{equation}

\paragraph{Rows $2N_f+1\ldots 4N_f$ --- vanishing ac current.}
A pure dc current bias forbids any ac current; every even harmonic $s=2h$ with
$h=1,\dots,N_f$ of \eqref{eq:Is} must vanish:
\begin{equation}
  \boxed{\;I_{2h}(\Wm)=0\;}\qquad h=1,\dots,N_f,
  \qquad
  \begin{aligned}
    F_{2N_f+h} &= \Real I_{2h},\\
    F_{3N_f+h} &= \Imag I_{2h}.
  \end{aligned}
  \label{eq:F_current}
\end{equation}

Equations \eqref{eq:F_reC}--\eqref{eq:F_current} are $\,(2N_f-1)+1+2N_f=4N_f\,$
real equations for the $4N_f$ real unknowns \eqref{eq:unknowns}: a square system.
The dc current $I_0$ is \emph{not} constrained --- it is the output read off the
converged solution, and sweeping the bias $\Om=\langle V\rangle$ traces the
$I$--$V$ curve.

\section{Newton solution and the analytic Jacobian}
\label{sec:jacobian}

\subsection{The solver}

The square system $\mathbf F(\mathbf x)=\mathbf 0$ is solved by
\texttt{NLsolve.nlsolve} with \texttt{method=:trust\_region}, supplying both
$\mathbf F$ (\texttt{IbiasResidual\_Tfull}) and its exact Jacobian
$J_{ij}=\partial F_i/\partial x_j$ (\texttt{IbiasJacobian\_Tfull}).  A trust-region Newton
step solves
\begin{equation}
  J(\mathbf x_k)\,\delta\mathbf x = -\,\mathbf F(\mathbf x_k),
  \qquad
  \mathbf x_{k+1} = \mathbf x_k + \delta\mathbf x
\end{equation}
within an adaptively sized trust region (so divergent Newton steps near the
strongly nonlinear gap edges are damped).  Convergence uses tolerances
\texttt{ftol}, \texttt{xtol}.  The bias is swept from high to low
$\langle V\rangle$ and the converged $\mathbf x$ at one bias seeds the next solve
(numerical continuation), which keeps the iterate inside the basin of attraction
across the subharmonic features.

\subsection{Jacobian of the algebraic constraints}

Rows \eqref{eq:F_reC}--\eqref{eq:F_gauge} are polynomial in $\mathbf x$, so their
derivatives are written in closed form.  Differentiating \eqref{eq:Cdef} and
using $\Wm_{r_q}=x_q+\ii x_{2N_f+q}$:
\begin{align}
  \frac{\partial F_h}{\partial \Real \Wm_q} &= \Real \Wm_{q+2h}+\Real \Wm_{q-2h},
  & \frac{\partial F_h}{\partial \Imag \Wm_q} &= \Imag \Wm_{q+2h}+\Imag \Wm_{q-2h},
  \label{eq:jac_reC}\\
  \frac{\partial F_{N_f-1+h}}{\partial \Real \Wm_q} &= \Imag \Wm_{q+2h}-\Imag \Wm_{q-2h},
  & \frac{\partial F_{N_f-1+h}}{\partial \Imag \Wm_q} &= -\Real \Wm_{q+2h}+\Real \Wm_{q-2h},
  \label{eq:jac_imC}\\
  \frac{\partial F_{2N_f-1}}{\partial \Real \Wm_q} &= 2\,\Real \Wm_q,
  & \frac{\partial F_{2N_f-1}}{\partial \Imag \Wm_q} &= 2\,\Imag \Wm_q,
  \label{eq:jac_norm}\\
  \frac{\partial F_{2N_f}}{\partial \Real \Wm_q} &= 0,
  & \frac{\partial F_{2N_f}}{\partial \Imag \Wm_q} &= 1,
  \label{eq:jac_gauge}
\end{align}
(terms with out-of-range harmonic index are dropped).  These are exactly the
hand-coded blocks at the top of \texttt{IbiasJacobian\_Tfull}.

\subsection{Jacobian of the current rows}

The nontrivial part is $\partial I_{2h}/\partial x_j$.  From \eqref{eq:Ihat},
\begin{equation}
  \frac{\partial \hat I}{\partial x_j}
  = -\,\frac{\partial \check V_i}{\partial x_j}\,\check G^<
    \;-\;\check V_i\,\frac{\partial \check G^<}{\partial x_j}.
  \label{eq:dIhat}
\end{equation}
Two simplifications make this cheap and \emph{exact}:

\paragraph{(a) $\check V$ and $\check V_i$ are linear in $\mathbf x$.}
Hence
\begin{equation}
  M^{(i)}_j \equiv \frac{\partial \check V_i}{\partial x_j},
  \qquad
  M_j \equiv \frac{\partial \check V}{\partial x_j}
  \label{eq:Mdef}
\end{equation}
are \emph{constant}, $\mathbf x$-independent, and \emph{sparse}: each unknown
$x_j$ (one real or imaginary part of one $\Wm_r$) sits on a single pair of
super/sub-diagonals of the $L/R$-off-diagonal blocks of \eqref{eq:hop_floquet}.
The code obtains them once, by a forward finite difference that is exact for a
linear map,
\begin{equation}
  M_j = \frac{\check V(\mathbf x+\hat e_j)-\check V(\mathbf x)}{\delta V}\Big|_{\delta V=1},
  \qquad
  M^{(i)}_j = \frac{\check V_i(\mathbf x+\hat e_j)-\check V_i(\mathbf x)}{\delta V}\Big|_{\delta V=1},
\end{equation}
(arrays \texttt{Mar}, \texttt{Miar}).  The thresholding
\verb|.* (abs.(Mar) .> 0.5*T*deltaV)| merely discards numerical-zero entries: the
genuine nonzero entries all have magnitude $\mathcal T$.  \texttt{Miar} is stored
with an overall minus sign ($\texttt{Miar}=-M^{(i)}_j$) to fold the leading minus
of \eqref{eq:dIhat} into the matrix.

\paragraph{(b) The lesser self-energy does not depend on $\mathbf x$.}
From \eqref{eq:sigl_bath}--\eqref{eq:sigl_surf}, $\check\Sigma^<$ is built solely
from the bare lead propagators and $\Gamma$; it contains no $\Wm$.  Therefore
$\partial\check\Sigma^</\partial x_j=0$.  Differentiating the Dyson equation
\eqref{eq:dyson} (with $\check G^r=(\id-\check g^r\check V)^{-1}\check g^r$,
so $\partial\check G^r=\check G^r M_j\check G^r$) and the Keldysh equation
\eqref{eq:keldysh}, and noting that the same hopping matrix $\check V$ enters the
advanced sector ($\partial\check G^a=\check G^a M_j\check G^a$), gives the compact
result
\begin{equation}
  \boxed{\;
  \frac{\partial \check G^<}{\partial x_j}
  = \check G^r M_j\,\check G^< \;+\; \check G^<\,M_j\,\check G^a \;}
  \label{eq:dGl}
\end{equation}
(the cross terms recombine using $\check G^<=\check G^r\check\Sigma^<\check G^a$).
Substituting \eqref{eq:dGl} into \eqref{eq:dIhat},
\begin{equation}
  \frac{\partial \hat I}{\partial x_j}
  = -\,M^{(i)}_j\,\check G^<
    \;-\;\check V_i\Big(\check G^r M_j\,\check G^< + \check G^<M_j\,\check G^a\Big).
  \label{eq:dIhat_full}
\end{equation}
This is the exact derivative.  Equation \eqref{eq:dIhat_full} is the commented
``\texttt{m0}'' line in \texttt{IbiasJacobian\_Tfull}.  The line that is actually executed
(``\texttt{m2}'') is the algebraically identical reordering
\begin{equation}
  \frac{\partial \hat I}{\partial x_j}
  = \big(\underbrace{-\check V_i\check G^<}_{\texttt{VviGlwmn}}\big)M_j\,\check G^a
   +\big(\underbrace{-\check V_i\check G^r}_{\texttt{VviGrwmn}}\big)M_j\,\check G^<
   +\big(\underbrace{-M^{(i)}_j}_{\texttt{Miar}}\big)\check G^<,
  \label{eq:dIhat_m2}
\end{equation}
in which the two factors $-\check V_i\check G^r$ and $-\check V_i\check G^<$ are
formed \emph{once} per energy $\omega$ (outside the loop over the $4N_f$
variables $j$), leaving only the cheap multiplications by the sparse $M_j$ inside
the inner loop.  This is the key performance optimization of \texttt{IbiasJacobian\_Tfull}.

\paragraph{Assembling the current-row Jacobian.}
For each variable $j$, the derivative matrix \eqref{eq:dIhat_m2} is reduced to
Floquet-mode current derivatives exactly as in
\eqref{eq:Imn_omega}--\eqref{eq:Is}: take the Nambu trace of the $LL$ minus $RR$
block per Floquet pair $(m,n)$ (array \texttt{jacIiwf}), integrate over $\omega$
(array \texttt{jacIif}), and sum over $n-m=2h$ (array \texttt{jacIfa}).  Finally
\begin{equation}
  J_{2N_f+h,\,j} = \Real\frac{\partial I_{2h}}{\partial x_j},
  \qquad
  J_{3N_f+h,\,j} = \Imag\frac{\partial I_{2h}}{\partial x_j},
  \qquad h=1,\dots,N_f.
  \label{eq:jac_current}
\end{equation}

\subsection{Perturbative Jacobians}
\label{sec:pert_jacobian}

For the truncated transparency schemes ($\texttt{ws}=2,4,6,8$;
\S\ref{sec:perturbation}) the exact resummed $\check G^<$ in
\eqref{eq:dGl}--\eqref{eq:dIhat_m2} is replaced by its finite expansion in powers
of $\check V$, built entirely from the \emph{bare} lead propagators
$\check g^r,\check g^a,\check g^<$.  Because the truncated current is then a finite
polynomial in $\mathbf x$ (the only $\mathbf x$-dependence is through the linear
$\check V$, $\check V_i$), its Jacobian is again \emph{exact} and is obtained by
the elementary product rule.  This subsection writes out that derivative and the
numerical scheme used by \texttt{IbiasJacobian\_T2/\_T4/\_T6/\_T8}.

\paragraph{Truncated lesser Green's function.}
Expanding the Dyson/Keldysh equations to finite order (\eqref{eq:pert_expansion})
groups the lesser GF by transparency order,
\begin{equation}
  \check G^<_{(\le 2P)} \;=\; \sum_{p=1}^{P}\check G^<_{(2p)},
  \qquad
  \check G^<_{(2p)} \;=\; \sum_{k=0}^{2p-1}
  (\check g^r\check V)^{k}\,\check g^<\,(\check V\check g^a)^{\,2p-1-k},
  \label{eq:Gl_trunc}
\end{equation}
with $P=1,2,3,4$ for $\texttt{ws}=2,4,6,8$.  Equation \eqref{eq:Gl_trunc} is what
\texttt{Glesser\_Floquet\_T2/\_T4/\dots} assemble; e.g.\ at $\mathcal T^2$,
$\check G^<_{(2)}=\check g^r\check V\check g^<+\check g^<\check V\check g^a$, and at
$\mathcal T^4$ the four chains $k=0,\dots,3$ of $\check G^<_{(4)}$ are added.

\paragraph{Derivative of the truncated current.}
The current operator is $\hat I=-\check V_i\,\check G^<_{(\le2P)}$, so by the
product rule
\begin{equation}
  \frac{\partial \hat I}{\partial x_j}
  = \underbrace{-\,M^{(i)}_j\,\check G^<_{(\le2P)}}_{\text{vertex term}}
   \;\underbrace{-\;\check V_i\,\frac{\partial \check G^<_{(\le2P)}}{\partial x_j}}_{\text{propagator term}},
  \label{eq:dIhat_pert}
\end{equation}
with $M^{(i)}_j=\partial\check V_i/\partial x_j$ (array \texttt{Miar}, sign-folded)
exactly as in \S\ref{sec:jacobian}.  Since the bare $\check g$'s carry no
$\mathbf x$, differentiating a single chain in \eqref{eq:Gl_trunc} replaces each
of its $2p-1$ factors of $\check V$ in turn by $M_j=\partial\check V/\partial x_j$
(array \texttt{Mar}):
\begin{equation}
  \frac{\partial}{\partial x_j}\Big[(\check g^r\check V)^{k}\check g^<(\check V\check g^a)^{2p-1-k}\Big]
  = \sum_{\text{each }\check V\text{ slot}} \underbrace{(\cdots)}_{L}\,M_j\,\underbrace{(\cdots)}_{R},
  \label{eq:chain_rule}
\end{equation}
i.e.\ a sum of terms $L\,M_j\,R$ in which $L$ is the sub-chain to the \emph{left}
of the replaced $\check V$ (ending in a bare $\check g$) and $R$ the sub-chain to
its \emph{right} (starting with a bare $\check g$).

\paragraph{Factorized (computed) form.}
Multiplying \eqref{eq:chain_rule} on the left by $-\check V_i$ and summing all
chains and slots, every contribution to the propagator term of
\eqref{eq:dIhat_pert} has the form
\begin{equation}
  \big(-\check V_i\,L\big)\;M_j\;R,
  \label{eq:LMR}
\end{equation}
where the left factors $-\check V_i L$ and the right factors $R$ are the
\emph{same handful of chains for every variable $j$}.  They are therefore built
\emph{once per frequency} and reused across the $4N_f$ columns; only the sparse
multiply by $M_j$ sits in the inner loop.

\paragraph{Order $\mathcal T^4$ written out.}
Take $\check G^<_{(4)}$, whose $k=3$ chain is $(\check g^r\check V)^3\check g^<$
(three $\check V$'s).  Replacing each $\check V$ and prepending $-\check V_i$ gives
\begin{equation}
  -\check V_i\,\partial_j\!\big[(\check g^r\check V)^3\check g^<\big]
  = \underbrace{(-\check V_i\check g^r)}_{\texttt{Vvigr}}M_j\underbrace{(\check g^r\check V)^2\check g^<}_{\texttt{gr\_s\_gr\_s\_gl}}
  + \underbrace{(-\check V_i\check g^r\check V\check g^r)}_{\texttt{Vvigr\_s\_gr}}M_j\underbrace{\check g^r\check V\check g^<}_{\texttt{gr\_s\_gl}}
  + \underbrace{(-\check V_i(\check g^r\check V)^2\check g^r)}_{\texttt{Vvigr\_s\_gr\_s\_gr}}M_j\underbrace{\check g^<}_{\texttt{glwmn}} ,
  \label{eq:T4_k3}
\end{equation}
which is exactly one line of the kernel in \texttt{IbiasJacobian\_T4}.  Here the
naming convention is: \texttt{s} denotes the hopping $\check V$ and
\texttt{r/l/a} denote $\check g^r/\check g^</\check g^a$, so
\texttt{gr\_s\_gl}$=\check g^r\check V\check g^<$ and
\texttt{Vvigr\_s\_gr}$=-\check V_i\,\check g^r\check V\check g^r$.  The other three
chains ($k=2,1,0$, ending in $\check g^a$ factors) contribute the remaining
lines; together with the $\mathcal T^2$ pair $-\check V_i\check g^r M_j\check g^<$,
$-\check V_i\check g^< M_j\check g^a$ and the vertex term $-M^{(i)}_j\check
G^<_{(4)}$ (code \texttt{Miar*Glwmn\_w4}) this is the full
$\partial\hat I/\partial x_j$ at $\mathcal T^4$ --- $15$ matrix products plus the
vertex term.  At $\mathcal T^6,\mathcal T^8$ the same construction adds the order
$\mathcal T^6,\mathcal T^8$ chains; the kernels group them as
\verb|t3_1| ($\mathcal T^2$), \verb|t3_3| ($\mathcal T^4$), \verb|t3_5| ($\mathcal
T^6$), \verb|t3_7| ($\mathcal T^8$) and \verb|t3_f| (vertex term).

\paragraph{Numerical manipulations.}
For a given bias the routine performs, for each frequency $\omega\in$
\texttt{war0}:
\begin{enumerate}
  \item build the bare propagators $\check g^r$ (\texttt{grwmn}),
        $\check g^a=\check g^{r\dagger}$ (\texttt{gawmn}), $\check g^<$
        (\texttt{glwmn}), and the truncated $\check G^<_{(\le2P)}$
        (\texttt{Glwmn\_wN} via \texttt{Glesser\_Floquet\_TN});
  \item build the right chains $R$ --- e.g.\ \texttt{gr\_s\_gl}, \texttt{gl\_s\_ga},
        \texttt{gr\_s\_gr\_s\_gl}, \dots --- by repeated multiplication by
        $\check V$ (\texttt{Vv});
  \item build the left chains $-\check V_i L$ --- \texttt{Vvigr},
        \texttt{Vvigr\_s\_gr}, \dots --- by one $\texttt{mul!}$ each
        (\emph{computed once}, reused for all $j$; at $\mathcal T^6/\mathcal T^8$
        the buffers are overwritten in place with \verb|.=|);
  \item loop (threaded) over the $4N_f$ variables $j$: assemble the derivative
        matrix as the order-grouped sum
        $\texttt{t3}=\texttt{t3\_1}+\texttt{t3\_3}+\dots$ of the terms
        \eqref{eq:LMR} with the sparse $M_j$ (\texttt{Mar[:,:,j]}) plus the vertex
        term \texttt{Miar[:,:,j]*Glwmn\_wN};
  \item reduce \texttt{t3} to the Floquet-mode current derivative by the
        $LL\!-\!RR$ Nambu trace and accumulate the frequency integral on the fly
        into \texttt{jacIif} (Tier-2 accumulation of \S\ref{sec:jacobian}).
\end{enumerate}
Finally \texttt{jacIif} is scaled by $\Delta\omega/2\pi$, collapsed onto current
harmonics (\texttt{jacIfa}, summing $n-m=2h$), and written to the current rows
\eqref{eq:jac_current}.

Two implementation notes.  (i) \texttt{IbiasJacobian\_T2} does not need the chain
machinery: $\partial_j\check G^<_{(2)}=\check g^r M_j\check g^<+\check g^< M_j\check
g^a$ has the same three-term shape as the exact case, so it reuses the
sparse-reordered form \eqref{eq:dIhat_m2} verbatim but with bare propagators and
$\check G^<_{(2)}$ in the vertex term.  (ii) \texttt{IbiasJacobian\_T6/\_T8} are
provided for completeness; \texttt{phisolve} currently runs $\texttt{ws}=6,8$
\emph{without} supplying an analytic Jacobian (\texttt{NLsolve} then finite-
differences the residual), and the $\mathcal T^6$ kernel omits the vertex term
that the $\mathcal T^8$ kernel keeps (\verb|t3_f|).

\section{Perturbative (transparency) expansion}
\label{sec:perturbation}

To dissect which microscopic processes produce each feature of the
$I$--$V$ curve, the lesser Green's function is also evaluated as a power series
in the transparency (a Werthamer-type tunnelling expansion).  Expanding the
Dyson equation \eqref{eq:dyson} and Keldysh equation \eqref{eq:keldysh} in
powers of $\Sig^r\propto\mathcal{T}$,
\begin{equation}
  \Gr^< = \sum_{p\ge 1}
  \sum_{k=0}^{2p-1}
  (\gr^r\Sig^r)^{k}\,\gr^<\,(\Sig^r\gr^a)^{2p-1-k},
  \label{eq:pert_expansion}
\end{equation}
and truncating at order $\mathcal{T}^{2}$, $\mathcal{T}^{4}$, $\mathcal{T}^{6}$,
$\mathcal{T}^{8}$ gives the functions \texttt{Glesser\_Floquet\_T2},
\texttt{\_T4}, \texttt{\_T6}, \texttt{\_T8} respectively (and the corresponding
currents \texttt{current\_Floquet\_T2}\,$\ldots$, self-consistency equations
\texttt{IbiasResidual\_T2}\,$\ldots$).  The order is selected by the \texttt{ws} flag
(\texttt{ws}\,$=0$ is the exact Dyson resummation; \texttt{ws}\,$=2,4,6,8$ the
truncated orders).  The lowest nontrivial order $\mathcal{T}^2$ is the
single-quasiparticle (and lowest pair) contribution; the $\mathcal{T}^4$
order is where the \emph{two-quasiparticle} processes that the paper identifies
as the origin of the current-biased subharmonic gap structure first appear.

\section{Observables and the subharmonic gap structure}
\label{sec:observables}

The principal outputs are:
\begin{itemize}
  \item the dc $I$--$V$ curve $I_{\rm dc}(\langle V\rangle)$, usually plotted in
        the dimensionless form $I\,e R_N/\Delta$ versus $eV/\Delta$;
  \item the differential conductance $\dd I/\dd V$, in which the subharmonic gap
        structure (SGS) appears as a series of features at
        \begin{equation}
          eV = \frac{2\Delta}{n},\qquad n=1,2,3,\dots
        \end{equation}
        (the dashed vertical lines in the plotting code).
\end{itemize}
The normal-state resistance $R_N$ used for normalization is obtained from the
same machinery by setting $\Delta=0$ and reading off the linear (Ohmic)
response, $R_N = \langle V\rangle / I_{\rm dc}\big|_{\Delta=0}$ (routine
\texttt{current\_full\_RN}).

\section{Extension to the \texorpdfstring{$4\times4$}{4x4} Nambu\texorpdfstring{$\otimes$}{x}spin basis: classical-spin (YSR) impurities}
\label{sec:ysr}

All of the preceding development uses a $2\times2$ Nambu space, which is
sufficient for a clean spin-singlet junction.  To host \emph{classical-spin
(magnetic) impurities} --- and thereby Yu--Shiba--Rusinov (YSR) bound states and
the Josephson diode effect --- the spin degree of freedom must be made explicit.
This section documents the $4\times4$ Nambu$\otimes$spin extension implemented in
\texttt{Keldyshsetup\_Floquetn\_ext.jl}, following Eqs.~(7)--(8) of
\href{https://arxiv.org/abs/2602.15213}{arXiv:2602.15213}.  The structure of
\S\S\ref{sec:gf}--\ref{sec:jacobian} is reused verbatim; we highlight only what
changes.

\subsection{Basis and the spin-resolved surface Green's function}
\label{sec:ysr_basis}

The Nambu spinor $\Psi=(c_\uparrow,c^\dagger_\downarrow)^{\mathsf T}$ of
\S\ref{sec:setup} is replaced by the full Nambu$\otimes$spin spinor
\begin{equation}
  \check c \;=\; \big(c_\uparrow,\,c_\downarrow,\,c^\dagger_\uparrow,\,c^\dagger_\downarrow\big)^{\mathsf T}
  \;=\;\tau\otimes\sigma,
\end{equation}
with $\tau$ the particle--hole and $\sigma$ the spin structure (particle block
$(\uparrow,\downarrow)$, hole block $(\uparrow^\dagger,\downarrow^\dagger)$).  The
clean BCS surface Green's function \eqref{eq:surface_gf} becomes the $4\times4$
matrix
\begin{equation}
  \check g_0(z) \;=\; \frac{1}{\zeta\sqrt{\Delta^2-z^2}}
  \begin{pmatrix}
    -z & 0 & 0 & \Delta\\
    0 & -z & -\Delta & 0\\
    0 & -\Delta & -z & 0\\
    \Delta & 0 & 0 & -z
  \end{pmatrix},
  \qquad z=\omega+\ii\Gamma,
  \label{eq:g0_4x4}
\end{equation}
in which the singlet pairing couples $c_\uparrow\!\leftrightarrow\!c^\dagger_\downarrow$
with $+\Delta$ (\emph{block~A}) and $c_\downarrow\!\leftrightarrow\!c^\dagger_\uparrow$
with $-\Delta$ (\emph{block~B}).  Equation~\eqref{eq:g0_4x4} is just the two
time-reversed copies of \eqref{eq:surface_gf} embedded with opposite-sign gaps;
it is the helper \texttt{gsurf4} evaluated at $\check V_{\rm imp}=0$.

\subsection{Classical-spin impurity and the local Dyson dressing}
\label{sec:ysr_imp}

A classical impurity spin (exchange $\mathbf J=(J_x,J_y,J_z)$, potential $K$)
enters as a static, local self-energy [Eq.~(7) of arXiv:2602.15213]
\begin{equation}
  \check V_{\rm imp}
  \;=\; \sum_{k=x,y,z} J_k\,\big(\tau_u\otimes\sigma_k-\tau_d\otimes\sigma_k^{*}\big)
        \;+\; K\,\big(\tau_z\otimes\sigma_0\big)
  \;=\;
  \begin{pmatrix}
    \mathbf J\!\cdot\!\boldsymbol\sigma & 0\\
    0 & -(\mathbf J\!\cdot\!\boldsymbol\sigma)^{*}
  \end{pmatrix}
  + K
  \begin{pmatrix}\sigma_0 & 0\\ 0 & -\sigma_0\end{pmatrix},
  \label{eq:Vimp}
\end{equation}
with the particle/hole projectors $\tau_{u}=\half(\tau_0+\tau_z)$,
$\tau_{d}=\half(\tau_0-\tau_z)$ (helper \texttt{impurity4}).  The longitudinal
part $J_z$ is diagonal in spin and keeps blocks A and B decoupled (it only shifts
them oppositely, splitting the YSR levels); the transverse parts $J_x,J_y$ are
off-diagonal in spin and \emph{couple} A and B.

The impurity is resummed into the surface propagator by a local Dyson (T-matrix)
dressing,
\begin{equation}
  \boxed{\;\check g(z) \;=\; \big(\id_4-\check g_0(z)\,\check V_{\rm imp}\big)^{-1}\check g_0(z)\;}
  \label{eq:local_dyson}
\end{equation}
(the body of \texttt{gsurf4}), whose in-gap poles are the YSR bound states.  Each
lead is dressed \emph{independently}: the left lead uses $(\mathbf J_L,K_L)$ and
the right lead $(\mathbf J_R,K_R)$, so the two contacts are in general
inequivalent (code arguments \texttt{JL,KL,JR,KR}).

\subsection{Floquet / lead layout}
\label{sec:ysr_floquet}

The Floquet construction of \S\ref{sec:floquet} is unchanged except for the
larger inner block.  The bare lead propagators are block-diagonal in the Floquet
index,
\begin{equation}
  [\check g^r]_{mn}=\delta_{mn}\,\check g(\omega+m\Om+\ii\Gamma),
\end{equation}
with $\check g$ from \eqref{eq:local_dyson}, and the lesser function
$\check g^<=-(\check g^r-\check g^a)\,\theta\!\big(-(\omega+m\Om)\big)$ as in
\eqref{eq:bare_lesser} --- now automatically including the equilibrium filling of
any YSR level.  The working space grows to
\begin{equation}
  \underbrace{2}_{L,R}\times\underbrace{4}_{\text{Nambu}\otimes\text{spin}}\times(2N_f+1)
  \;=\;8(2N_f+1),
\end{equation}
versus $4(2N_f+1)$ in the $2\times2$ theory.  Routines: \texttt{grwmnf},
\texttt{gawmnf}, \texttt{glwmnf} (now taking \texttt{JL,KL,JR,KR}).

\subsection{Hopping and current vertex are spin-conserving}
\label{sec:ysr_hop}

The inter-lead tunnelling [Eq.~(8) of arXiv:2602.15213] conserves spin,
$\mathcal T\,(W\,\tau_u-W^{*}\tau_d)\otimes\sigma_0$, so the Floquet hopping
self-energy \eqref{eq:hop_floquet} is simply the $2\times2$ Nambu matrix tensored
with the spin identity,
\begin{equation}
  \big[\check V\big]^{LR}_{mn}
  =\mathcal T\begin{pmatrix}\Wm_{\,n-m}&0\\0&-\Wm^{*}_{\,m-n}\end{pmatrix}\!\otimes\sigma_0,
  \qquad
  \big[\check V_i\big]^{LR}_{mn}
  =\mathcal T\begin{pmatrix}\Wm_{\,n-m}&0\\0&+\Wm^{*}_{\,m-n}\end{pmatrix}\!\otimes\sigma_0,
  \label{eq:hop_4x4}
\end{equation}
and likewise for the $RL$ blocks.  In particular $\check V$ and $\check V_i$ carry
\emph{no} impurity parameters (these live only in the leads), so
\texttt{Vwmnf}/\texttt{Viwmnf} keep their $2\times2$-era signatures, differing
only by the \verb|kron(.,sig0)|.  They remain \emph{linear} in $\{\Wm_r\}$, which
is what preserves the entire Jacobian construction below.

\subsection{Dyson, Keldysh, and the current}
\label{sec:ysr_current}

The Dyson equation \eqref{eq:dyson}, the lesser self-energy
\eqref{eq:sigl_bath}--\eqref{eq:sigl_surf} (with $\tau_0\!\to\!\id_4$ and
$\tauz\!\to\!\tau_z\otimes\sigma_0$), and the Keldysh equation \eqref{eq:keldysh}
are formally identical, now solved in the $8(2N_f+1)$-dimensional space
(\texttt{Grwmnf}, \texttt{Siglf}, \texttt{Glesser\_Floquet\_Tfull}).  The current
\eqref{eq:current_time} is unchanged in form, $\hat I=-\check V_i\check G^<$, but
the Nambu trace becomes a Nambu$\otimes$spin trace over \emph{four} diagonal
components, with the right lead at offset $4(2N_f+1)$:
\begin{equation}
  I_{mn}(\omega)
  = \sum_{a=1}^{4}\big[\hat I^{LL}_{mn}\big]_{aa}
  - \sum_{a=1}^{4}\big[\hat I^{RR}_{mn}\big]_{aa},
  \label{eq:Imn_4x4}
\end{equation}
to be compared with the two-component trace \eqref{eq:Imn_omega}.  Everything
downstream --- the frequency integral \eqref{eq:Imn}, the harmonic collapse
\eqref{eq:Is}, the dc current --- is identical (routine
\texttt{current\_Floquet\_Tfull}, and the perturbative
\texttt{current\_Floquet\_T2/\_T4}).

\subsection{Self-consistency and Jacobian: what changes and what does not}
\label{sec:ysr_jac}

The phase $\phi(t)$ is a \emph{scalar} (spin-independent) gauge field, so the
unknown vector \eqref{eq:unknowns} is \emph{unchanged}: the same $4N_f$ reals.
Consequently, in the residual $\mathbf F(\mathbf x)$ of \S\ref{sec:explicit},
\begin{itemize}
  \item the unitarity rows \eqref{eq:F_reC}--\eqref{eq:F_norm} and the gauge row
        \eqref{eq:F_gauge} are \emph{identical} (they involve only $\{\Wm_r\}$,
        not the basis);
  \item only the ac-current rows \eqref{eq:F_current} change --- through the
        $4\times4$ current \eqref{eq:Imn_4x4} and the per-lead impurity.
\end{itemize}
The analytic Jacobian inherits the same split.  The constraint-row derivatives
\eqref{eq:jac_reC}--\eqref{eq:jac_gauge} are unchanged.  For the current rows the
two key facts of \S\ref{sec:jacobian} still hold verbatim: $\check\Sigma^<$ is
independent of $\mathbf x$, and $\check V,\check V_i$ are linear in $\mathbf x$,
so
\begin{equation}
  \frac{\partial\check G^<}{\partial x_j}=\check G^r M_j\check G^<+\check G^< M_j\check G^a,
  \qquad
  \frac{\partial\hat I}{\partial x_j}
  =(-\check V_i\check G^<)M_j\check G^a+(-\check V_i\check G^r)M_j\check G^<+(-M^{(i)}_j)\check G^<,
  \label{eq:dG_4x4}
\end{equation}
exactly Eqs.~\eqref{eq:dGl} and \eqref{eq:dIhat_m2}.  The \emph{only}
modifications are dimensional: the constant sparse derivatives $M_j,M^{(i)}_j$ and
all scratch matrices grow from $4(2N_f+1)$ to $8(2N_f+1)$, and the
per-Floquet-pair reduction uses the four-component Nambu$\otimes$spin trace
\eqref{eq:Imn_4x4}.  The perturbative chain derivatives of
\S\ref{sec:pert_jacobian} carry over identically (wider matrices, four-component
trace).  Routines: \texttt{IbiasResidual\_Tfull}, \texttt{IbiasJacobian\_Tfull},
and the truncated \texttt{\_T2/\_T4}; \texttt{phisolve} dispatches
$\texttt{ws}=0,2,4$ with \texttt{JL,KL,JR,KR} threaded through.

\subsection{Reduction to the \texorpdfstring{$2\times2$}{2x2} theory, and the diode effect}
\label{sec:ysr_check}

For a non-magnetic contact ($\mathbf J_L=\mathbf J_R=0$, $K_L=K_R=0$) the dressing
\eqref{eq:local_dyson} is trivial and $\check g_0$ \eqref{eq:g0_4x4} is
block-diagonal into block A --- which is \emph{precisely} the original $2\times2$
problem \eqref{eq:surface_gf} --- and its time-reverse block B.  Because the
hopping \eqref{eq:hop_4x4} is $\propto\sigma_0$, the two blocks never mix, and the
four-component trace \eqref{eq:Imn_4x4} equals twice the single-block trace.
Hence every observable computed in the extended module reduces to \emph{exactly
twice} its $2\times2$ counterpart --- the spin-degeneracy factor.  This is the
regression used to validate the port: with $\mathbf J=K=0$,
\texttt{current\_Floquet\_T2/\_T4} and the current rows of
\texttt{IbiasResidual}/\texttt{IbiasJacobian} reproduce $2\times$ the
\texttt{Keldyshsetup\_Floquetn.jl} values to machine precision, and the analytic
Jacobian matches a finite difference of the residual.

The physics that genuinely requires the $4\times4$ machinery is the
\emph{non-reciprocal} (diode) regime: for non-collinear or unequal leads
($\mathbf J_L\nparallel\mathbf J_R$) blocks A and B are coupled and the left/right
symmetry is broken, $I_{LR}\neq-I_{RL}$, so the supercurrent is no longer odd in
the phase/voltage --- the Josephson diode effect that motivates the extension.

\subsection{Correspondence to the extended code}
\label{sec:ysr_codemap}

\begin{center}
\renewcommand{\arraystretch}{1.25}
\begin{tabular}{ll}
\hline
\textbf{Equation / object} & \textbf{Routine in \texttt{Keldyshsetup\_Floquetn\_ext.jl}}\\
\hline
Clean $4\times4$ surface GF \eqref{eq:g0_4x4} & \texttt{gsurf4} (with $\check V_{\rm imp}=0$)\\
Impurity self-energy $\check V_{\rm imp}$ \eqref{eq:Vimp} & \texttt{impurity4(J,K)}\\
Local Dyson dressing \eqref{eq:local_dyson} & \texttt{gsurf4(z,zeta,delta,Vimp)}\\
Dressed lead GFs $g^{r,a,<}$ & \texttt{grwmnf}, \texttt{gawmnf}, \texttt{glwmnf} (\texttt{...,JL,KL,JR,KR})\\
Hopping / current vertex \eqref{eq:hop_4x4} & \texttt{Vwmnf}, \texttt{Viwmnf} (\verb|kron(.,sig0)|)\\
Dyson / lesser SE / Keldysh & \texttt{Grwmnf}, \texttt{Siglf}, \texttt{Glesser\_Floquet\_Tfull}\\
Current, four-component trace \eqref{eq:Imn_4x4} & \texttt{current\_Floquet\_Tfull}\\
Transparency expansion ($\mathcal T^2,\mathcal T^4$) & \texttt{Glesser/current\_Floquet/IbiasResidual/IbiasJacobian\_T2,\_T4}\\
Self-consistency / Jacobian \eqref{eq:dG_4x4} & \texttt{IbiasResidual\_Tfull}, \texttt{IbiasJacobian\_Tfull}\\
Bias sweep + Newton solve ($\texttt{ws}=0,2,4$) & \texttt{phisolve(...,JL,KL,JR,KR)}\\
Normal-state resistance $R_N$ & \texttt{RN\_full(...,JL,KL,JR,KR)}\\
\hline
\end{tabular}
\end{center}

\paragraph{What differs from the $2\times2$ routines.}
The signatures gain the per-lead impurity arguments \texttt{JL,KL,JR,KR}; matrices
are sized $8(2N_f+1)$; the Nambu trace sums four diagonal components with
right-lead offset $4(2N_f+1)$; and \texttt{Vwmnf}/\texttt{Viwmnf} tensor the
$2\times2$ hopping with $\sigma_0$.  All other differential and algebraic
structure is identical to \S\S\ref{sec:explicit}--\ref{sec:jacobian}.

\section{Correspondence to the code}
\label{sec:codemap}

\begin{center}
\renewcommand{\arraystretch}{1.25}
\begin{tabular}{ll}
\hline
\textbf{Equation / object} & \textbf{Routine in \texttt{Keldyshsetup\_Floquetn.jl}}\\
\hline
Surface GF $g^r$, $g^a$ \eqref{eq:surface_gf} & \texttt{grwmnf}, \texttt{gawmnf}\\
Surface GF (transfer matrix) & \texttt{surfacegr}\\
Bare lesser $g^<$ \eqref{eq:bare_lesser} & \texttt{glwmnf}\\
Hopping self-energy $\Sigma_{LR/RL}$ \eqref{eq:hop_floquet} & \texttt{Vwmnf}, \texttt{Viwmnf}\\
Dyson eq. $\Gr^r$ \eqref{eq:dyson} & \texttt{Grwmnf}, \texttt{Gawmnf}\\
Lesser self-energy \eqref{eq:sigl_bath}, \eqref{eq:sigl_surf} & \texttt{Siglf}\\
Keldysh eq. $\Gr^<$ \eqref{eq:keldysh} & \texttt{Glesser\_Floquet\_Tfull}\\
Transparency expansion \eqref{eq:pert_expansion} & \texttt{Glesser\_Floquet\_T2/\_T4/\_T6/\_T8}\\
Current \eqref{eq:current_time}, \eqref{eq:current_harmonics} & \texttt{current\_Floquet\_Tfull}, \texttt{current\_Floquet\_T2/\ldots}\\
Self-consistency $\mathbf F(\mathbf x)$ \eqref{eq:unitarity}--\eqref{eq:noac} & \texttt{IbiasResidual\_Tfull}, \texttt{IbiasResidual\_T2/\ldots}\\
Jacobian of $\mathbf F$ & \texttt{IbiasJacobian\_Tfull}, \texttt{IbiasJacobian\_T2/\ldots}\\
Bias sweep + Newton solve & \texttt{phisolve}\\
Voltage waveform $V(t)$ & \texttt{Vt}\\
Normal-state resistance $R_N$ & \texttt{current\_full\_RN}\\
\hline
\end{tabular}
\end{center}

\paragraph{Index conventions.}
Matrices act on the ordered product space (lead)$\otimes$(Floquet)$\otimes$(Nambu).
A Floquet index $m\in\{-N_f,\dots,N_f\}$ maps to the block row/column
$(-m+N_f+1)$; the complex coefficient array \texttt{Vip} of length $4N_f+1$ is
centered at index $2N_f+1$ ($\leftrightarrow m=0$) and is nonzero only on odd
$m$.  The solver's real unknown vector \texttt{Vipi}/\texttt{Vipsol} stacks the
$2N_f$ real parts followed by the $2N_f$ imaginary parts of these odd
harmonics.

\section*{References}
\begin{enumerate}
  \item \emph{Origin of Subharmonic Gap Structure of DC Current-Biased
        Josephson Junctions},
        \href{https://arxiv.org/abs/2412.09862}{arXiv:2412.09862}.
  \item Surface Green's function / transfer-matrix construction:
        Phys.\ Rev.\ B \textbf{83}, 085412 (2011).
\end{enumerate}

\end{document}
